\section{Methodology}\label{sec:evaluation-methodology}

The evaluation compares two tagged compiler releases on identical synthetic workloads.
Version~1.0 corresponds to git tag \texttt{v1.0.0-base} and the \texttt{dslrc} CLI over \texttt{.dsl} sources;
version~2.0 corresponds to \texttt{v2.0.0-explicit-types} and the \texttt{motivoc} CLI over \texttt{.motivo} sources.
Both binaries were built as separate Release configurations from the same codebase and benchmarked on one machine without
switching versions between runs.

Measurements were taken on an Apple M3 Pro (arm64, macOS~26.5) with 18~GB of RAM.
Wall-clock times were recorded by a small benchmark harness that wraps each compiler
invocation in Python's \texttt{time.perf\_counter}, performs 5 warmup runs, and averages 50 measured runs.
Each measured command compiles one source file and writes a \texttt{.mid} file beside the input, so the timer includes
process startup, parsing, semantic analysis, lowering, MIDI serialization, and file I/O.

Per-stage breakdowns use in-process timers enabled through a compile-time benchmark flag in the compiler driver.
These timers report parse, semantic, lowering, and MIDI-write means independently of subprocess overhead; their sum is
therefore smaller than the end-to-end wall-clock total.

The fixtures are synthetic stress programs: they exercise expansion and throughput at the implementation's
20,000-event cap (Section~\ref{sec:limits}) rather than modeling finished compositions.
The evaluation therefore targets engineering properties (expansion ratio, compile time, and safety limits), not musical
quality or user preference.

\subsection{Theoretical Complexity}\label{detail-subsec:theoretical-complexity}

The compilation pipeline is composed of several sequential passes, each with a distinct computational complexity:

\topic{Frontend and Semantic Analysis.} The complexity of parsing and type-checking is proportional to the
number of lines of source code, $O(S)$, where $S$ is the size of the abstract syntax tree.
Because these passes do not unroll loops or instantiate patterns multiple times, they remain extremely fast regardless
of the composition's length.
Version~2.0 adds explicit type annotations and signature-based pattern resolution, but the traversal still visits each
AST node a constant number of times; the semantic pass was simplified from two traversals in version~1.0 to one.

\topic{Lowering Engine.} The unrolling process has a time complexity of $O(E)$, where $E$ is the total number
of generated \texttt{NoteEvent}s.
A simple \texttt{loop (1000)} statement will execute its body 1,000 times during this phase, meaning the lowering time
is proportional to the density of the final timeline, not the brevity of the source code.

\topic{MIDI Backend.} The backend must sort the absolute events by their tick values before computing
delta-times.
This stable sort introduces a time complexity of $O(E \log E)$.
Serializing the events into the binary file requires a final linear pass, $O(E)$.

Therefore, the overall time complexity of the compilation process is bounded by the sorting phase in the backend:
$O(E \log E)$.
